% generated by roam.py handbook export -- не редактировать, править заметки
% 20260923123518-теорема_веиерштрасса_о_достижении_экстремума.org (body fdabd43a4994)

\begin{theorem}[{Вейерштрасса о достижении экстремума}]\label[theorem]{note:82C70263-336A-49F5-82DB-9D25C84D6A6D}
  Пусть $K \subset \R^n$~--- непустое компактное (замкнутое и ограниченное) множество, $f: K \rightarrow \R$~--- непрерывная функция.
  Тогда $f$ ограничена на $K$ и достигает на нем наименьшего и наибольшего значений:
  \[
  \exists x_{*}, x^{*} \in K:\quad f(x_{*}) \le f(x) \le f(x^{*})\quad \forall x \in K
  \]

  \begin{itemize}
    \item $K$ не замкнуто: $f(x) = x$ на $(0;1]$~--- $\inf f = 0$ не достигается;
    \item $K$ не ограничено: $f(x) = e^{-x}$ на $[0;+\infty)$~--- $\inf f = 0$ не достигается;
    \item $f$ разрывна: $f(x) = x$ при $x \in [0;1)$, $f(1) = 0$~--- $\sup f = 1$ не достигается.
  \end{itemize}
\end{theorem}

\begin{remark}
  Все условия существенны:
\end{remark}

\begin{remark}
  Типичное применение в оптимизации: если $f$ растет на бесконечности, то вне большого шара ее значения заведомо больше, чем в какой-то точке внутри, и минимум достаточно искать на шаре (см. \hyperref[note:51CC880D-BF19-43E2-9599-FB31A9A7367F]{единственность минимума сильно выпуклой функции}).
\end{remark}

\begin{proof}
  Докажем для минимума; для максимума применим то же рассуждение к $-f$.

  \proofpart{Ограниченность снизу}
  Пусть $f$ не ограничена снизу: тогда $\forall k \in \mathbb{N}\ \exists x_k \in K:\ f(x_k) < -k$.
  $K$ ограничено, поэтому по теореме Больцано~--- Вейерштрасса у $(x_k)$ есть сходящаяся подпоследовательность $x_{k_j} \rightarrow \bar{x}$, а так как $K$ замкнуто, $\bar{x} \in K$.
  По непрерывности $f(x_{k_j}) \rightarrow f(\bar{x})$~--- конечное число, но $f(x_{k_j}) < -k_j \rightarrow -\infty$. Противоречие.

  \proofpart{Достижение}
  Пусть $m = \inf_K f > -\infty$. По определению инфимума $\forall k\ \exists x_k \in K:\ m \le f(x_k) < m + \frac{1}{k}$.
  Как и выше, выделим подпоследовательность $x_{k_j} \rightarrow x_{*} \in K$. По непрерывности
  \[
  f(x_{*}) = \lim_{j \rightarrow \infty} f(x_{k_j}) = m
  \]
  то есть $x_{*}$~--- точка минимума $f$ на $K$.
\end{proof}
